\section{Conclusion: The Trial That Never Ends}

Strip the story to what the evidence best supports, and it is still barely
credible. A seventeen-year-old, orphaned by the executioner and stripped of
his inheritance, taught himself into the chair of the most learned church in
the world; sold his books to live on four obols a day; taught catechumens by
day whom he then accompanied to their executions; learned Hebrew; built, with
a patron's fortune and a staff of stenographers, the first critical edition
of the Bible and the first Christian library of systematic exegesis; wrote
the first Christian systematic investigation of doctrine and the greatest
ancient defense of the faith; formed bishops for half the East and talked
straying ones back into the Church without a single anathema; was deposed by
his own bishop on grounds his allies never accepted; was tortured with
deliberate care under Decius and would not yield; and died, in the peace of
the Church but on the wrong side of Alexandria's sentence, at Tyre, at
sixty-nine.

The rest of the story happened to his name, and it is the part this study
has been most careful to state exactly. Origen was not tried in 553 and
convicted of heresy in his person by the assembled ecumenical Church; nor is
it true that he was never condemned. An imperial edict with synodal and
papal assent anathematized him by name in 543; an ecumenical council a
decade later listed him among heretics already condemned, while a set of
anathemas drawn up around that council condemned, under his name, a system
substantially built by men born three centuries after him, out of a book
that survives mainly in a translation its translator admitted to correcting.
Whether that amounts to ``Origen anathematized by the Church'' is a question
on which learned and faithful writers have divided for four hundred
years---``Many learned writers believe so; an equal number deny that they
were condemned''---and the honest report is the divided one.

Perhaps the most exact epitaph is the one his own method supplies. Origen
taught that the apostles had deliberately left questions open, and that the
lover of wisdom honors God by exploring them ``in the manner rather of an
investigation and discussion, than in that of fixed and certain decision.''
His gravest speculations were offered under that rule; posterity broke the
rule in both directions, first hardening the hypotheses into a system,
then condemning the system under his name. The Church kept what was
certain---his prayer, his senses of Scripture, his defense of freedom, his
witness under torture---and it says so in its Catechism; the questions he
left open are, many of them, open still. The boy who wrote to his father,
``Take heed not to change your mind on our account,'' never changed his own:
about the Scriptures, about freedom, about martyrdom, about the goodness of
God. What the Church finally does with such a mind is, on the evidence of
seventeen centuries, itself one of the questions left to be investigated.
